\documentclass[10pt,a4paper]{article}
\usepackage[a4paper,top=20mm,bottom=18mm,left=18mm,right=18mm,headheight=14pt]{geometry}
\usepackage{fontspec}
\usepackage{unicode-math}
\setmainfont{Cambria}
\setmathfont{Cambria Math}
\setmonofont{Consolas}[Scale=0.82]
\usepackage{polyglossia}
\setdefaultlanguage{english}

\usepackage{array,tabularx,booktabs}
\usepackage{microtype}
\usepackage{enumitem}
\usepackage{fancyhdr}
\usepackage{lastpage}
\usepackage{titlesec}
\usepackage[hidelinks,unicode]{hyperref}

\setlength{\parindent}{0pt}
\setlength{\parskip}{4pt}
\setlength{\emergencystretch}{2em}
\setlist[enumerate]{leftmargin=*,itemsep=2pt,topsep=3pt}
\renewcommand{\arraystretch}{1.18}
\newcolumntype{L}[1]{>{\raggedright\arraybackslash}p{#1}}
\newcolumntype{Y}{>{\raggedright\arraybackslash}X}
\newcommand{\code}[1]{{\ttfamily\small #1}}

\titleformat{\section}{\large\bfseries}{\thesection}{0.65em}{}
\titlespacing*{\section}{0pt}{10pt}{4pt}
\titleformat{\subsection}{\normalsize\bfseries}{\thesubsection}{0.6em}{}
\titlespacing*{\subsection}{0pt}{8pt}{3pt}

\pagestyle{fancy}
\fancyhf{}
\fancyhead[L]{\footnotesize LLSN-TN-001}
\fancyhead[C]{\footnotesize Rev. A}
\fancyhead[R]{\footnotesize 15 Nov 2026}
\fancyfoot[C]{\footnotesize Page \thepage\ of \pageref{LastPage}}
\renewcommand{\headrulewidth}{0.4pt}
\renewcommand{\footrulewidth}{0pt}
\setlength{\headsep}{10pt}

\begin{document}

\begin{center}
{\LARGE\bfseries Landing-Site Suitability Score\par}
\vspace{4pt}
{\large Fixed-Ideal TOPSIS Model for Lunar Site Screening\par}
\vspace{10pt}
{\normalsize MmilDeum Project Team\par}
{\small NASA Space Apps Challenge\par}
{\small 15 November 2026\par}
\end{center}

\section*{Abstract}
This technical note defines a configurable, candidate-independent index for screening lunar landing sites using solar illumination, direct-to-Earth visibility, outage duration, and terrain slope. Five measured quantities are normalized and combined by fixed-ideal TOPSIS, while mission limits are evaluated separately as feasibility constraints. Scores are comparable across candidate sets because the ideal and anti-ideal anchors are fixed. The model supports early trade studies; it is not a flight-certification method or a probability of mission success. Mission-specific limits, weights, power margins, thermal performance, communications design, landing dynamics, navigation uncertainty, and local hazards require separate engineering analysis.

\section*{Nomenclature}
\small
\begin{tabularx}{\linewidth}{@{}L{0.10\linewidth}Y L{0.12\linewidth}Y@{}}
\toprule
\textbf{Symbol} & \textbf{Definition} & \textbf{Symbol} & \textbf{Definition}\\
\midrule
$p$ & Candidate site $(\varphi,\lambda)$ & $[t_0,t_1]$ & Evaluation interval\\
$T$ & Interval duration & $\Delta t$ & Sampling step\\
$H_p(a)$ & Terrain-horizon elevation at azimuth $a$ & $u_S,u_E$ & Sun/Earth visibility indicators\\
$I,V$ & Illuminated and Earth-visible fractions & $D,B$ & Longest darkness and Earth-blackout durations\\
$\theta_{95}$ & 95th-percentile footprint slope & $\tau_P,\tau_C$ & Outage reference time scales\\
$\theta_{\mathrm{ref}}$ & Slope reference scale & $z_i,w_i$ & Normalized benefits and weights\\
$d_+,d_-$ & Distances to ideal and anti-ideal & $S(p)$ & Suitability score\\
$\mathcal{F}(p)$ & Mission-feasibility result & $\theta_{\max},D_{\max},B_{\max}$ & Configured mission limits\\
\bottomrule
\end{tabularx}
\normalsize

\section{Inputs and observable quantities}
For site $p$, interval $[t_0,t_1]$, and step $\Delta t$, the evaluation pipeline uses the following sources.

\small
\begin{tabularx}{\linewidth}{@{}L{0.23\linewidth}Y L{0.20\linewidth}@{}}
\toprule
\textbf{Source} & \textbf{Retrieved or derived values} & \textbf{Use}\\
\midrule
JPL Horizons API \cite{horizons} & Solar and terrestrial azimuth $a_{S,E}(t)$ and elevation $e_{S,E}(t)$ at the lunar observer & Visibility geometry\\
LOLA elevation data \cite{smith2010} & Elevation raster $z(x,y)$ and slope samples over the landing footprint & Terrain horizon and slope\\
Backend processing & $H_p(a)$, visibility intervals, outage durations, and $\theta_{95}$ & Scoring inputs and provenance\\
\bottomrule
\end{tabularx}
\normalsize

Absolute elevation and geological class are retained as descriptors but do not enter the score. Elevation data instead support the terrain-horizon and footprint-slope calculations. The global LOLA product \code{MAS\_225M\_16} reports median absolute along-track slope at a 225 m baseline in degrees \cite{rosenburg2011,pgda}. It is a site descriptor and is excluded from the score and feasibility constraints.

\newpage

\section{Visibility metrics}
Let $H_p(a)$ denote the terrain-horizon elevation at azimuth $a$, including observer height and lunar curvature. Define
\begin{align}
u_S(t)&=\mathbb{1}\!\left[e_S(t)>H_p(a_S(t))\right],\\
u_E(t)&=\mathbb{1}\!\left[e_E(t)>H_p(a_E(t))\right].
\end{align}
The indicators equal one when the corresponding body clears the local terrain horizon. The availability fractions are
\begin{equation}
I=\frac{1}{T}\int_{t_0}^{t_1}u_S(t)\,dt,
\qquad
V=\frac{1}{T}\int_{t_0}^{t_1}u_E(t)\,dt.
\end{equation}
The quantity $D$ is the \emph{longest continuous interval} for which $u_S=0$; $B$ is defined analogously for $u_E=0$. Both are measured in hours. Numerical sampling approximates the integrals and interval boundaries.

\section{Metric normalization}
The five scoring inputs are mapped to dimensionless benefits $z_i\in[0,1]$:
\begin{align}
z_1 &= I, & z_2 &= \exp\!\left(-\frac{D}{\tau_P}\right),\\
z_3 &= V, & z_4 &= \exp\!\left(-\frac{B}{\tau_C}\right),\\
z_5 &= \max\!\left[0,1-\left(\frac{\theta_{95}}{\theta_{\mathrm{ref}}}\right)^2\right].
\end{align}
$\theta_{95}$ is evaluated over the selected landing footprint. The initial reference values are $\tau_P=24$ h, $\tau_C=120$ h, and $\theta_{\mathrm{ref}}=10^\circ$. These parameters shape the scoring curves and are replaced by mission-derived values when vehicle requirements are available.

\section{Fixed-ideal TOPSIS score}
Let the nonnegative weights satisfy $\sum_{i=1}^{5}w_i=1$. The initial weight vector is
\begin{equation}
\boldsymbol{w}=(0.25,\ 0.20,\ 0.15,\ 0.10,\ 0.30),
\qquad
\boldsymbol{z}=(z_1,\ldots,z_5).
\end{equation}
Illumination and darkness account for 45\% of the weight, Earth visibility and blackout for 25\%, and terrain slope for 30\%. Each site is compared with the fixed ideal $\boldsymbol{1}$ and anti-ideal $\boldsymbol{0}$:
\begin{align}
d_+(p)&=\sqrt{\sum_{i=1}^{5}w_i\,[1-z_i(p)]^2},\\
d_-(p)&=\sqrt{\sum_{i=1}^{5}w_i\,z_i(p)^2}.
\end{align}
The suitability score is
\begin{equation}
\boxed{S(p)=100\,\frac{d_-(p)}{d_+(p)+d_-(p)}}.
\label{eq:score}
\end{equation}
Thus $0\leq S\leq100$, and adding or removing a candidate does not alter an existing site's score.

\newpage

\section{Constraints and data quality}
The numerical score and mission feasibility are reported separately. For configured limits,
\begin{equation}
\mathcal{F}(p)=
(\theta_{95}\leq\theta_{\max})\land
(D\leq D_{\max})\land
(B\leq B_{\max}).
\end{equation}
Darkness and blackout limits are applied only when specified. A site that fails a limit remains in the result set with \code{feasible=false} and the applicable entries in \code{failed\_constraints}; feasible sites rank ahead of infeasible sites.

If DEM coverage, ephemeris coverage, or a complete slope footprint is unavailable, the result is \code{insufficient\_data}. Each result records the DEM resolution, time step, footprint dimensions, and source identifiers.

\section{Comparison of two sites}
Sites $A$ and $B$ are evaluated with identical mission dates, sample interval, footprint definition, weights, and normalization parameters. Report
\begin{equation}
\Delta S=S(A)-S(B).
\end{equation}
Site $A$ \emph{Pareto-dominates} site $B$ when
\begin{equation}
\bigl(\forall i,\ z_i(A)\geq z_i(B)\bigr)
\land
\bigl(\exists j,\ z_j(A)>z_j(B)\bigr).
\end{equation}
When neither site dominates, the preferred site depends on the selected trade-offs and feasibility status.

\section{Worked example}
Consider two hypothetical sites with complete data and no failed mission constraints.

\begin{center}
\small
\begin{tabular}{@{}lrr@{}}
\toprule
\textbf{Measurement} & \textbf{Site A} & \textbf{Site B}\\
\midrule
Illuminated fraction $I$ & $0.80$ & $0.72$\\
Longest darkness $D$ & $8$ h & $3$ h\\
Earth-visible fraction $V$ & $0.45$ & $0.63$\\
Longest Earth blackout $B$ & $120$ h & $80$ h\\
Footprint slope $\theta_{95}$ & $4^\circ$ & $6^\circ$\\
\midrule
Normalized $\boldsymbol{z}$ & $(.800,.717,.450,.368,.840)$ & $(.720,.882,.630,.513,.640)$\\
\bottomrule
\end{tabular}
\end{center}
Evaluation of Equation~\eqref{eq:score} gives
\begin{equation}
S(A)=67.60,\qquad S(B)=68.40,\qquad \Delta S=-0.80.
\end{equation}
Site B ranks slightly higher under the initial weights, and neither site Pareto-dominates the other. Site A has greater total illumination and gentler terrain; Site B has shorter outages and better Earth visibility.

\newpage

\appendix
\section{Software interface and verification}
The reference interface is \code{score\_site(metrics, mission\_config)}. Its result contains \code{score}, \code{feasible}, \code{failed\_constraints}, \code{normalized\_metrics}, \code{model\_version}, and \code{data\_quality}. The companion function \code{compare\_sites(A, B, config)} returns both independent scores, $\Delta S$, feasibility states, a preferred site, and Pareto dominance.

Minimum verification cases are:
\begin{enumerate}
\item $\boldsymbol{z}=\boldsymbol{0}\Rightarrow S=0$ and $\boldsymbol{z}=\boldsymbol{1}\Rightarrow S=100$.
\item Adding a candidate does not change an existing score.
\item Feasibility limits take precedence over the numerical score.
\item Increasing $D$ or $B$, with other inputs fixed, cannot improve the score.
\item Invalid fractions, durations, scales, weights, and incomplete source coverage are rejected or reported as insufficient data.
\item The worked example is reproduced within $0.01$ point.
\end{enumerate}

\section*{Revision history}
\small
\begin{tabularx}{\linewidth}{@{}L{0.13\linewidth}L{0.20\linewidth}Y L{0.20\linewidth}@{}}
\toprule
\textbf{Revision} & \textbf{Date} & \textbf{Description} & \textbf{Prepared by}\\
\midrule
A & 15 Nov 2026 & Initial controlled release of scoring model v1.0 & MmilDeum Project Team\\
\bottomrule
\end{tabularx}
\normalsize

\begin{thebibliography}{9}
\small
\bibitem{rosenburg2011}
M. A. Rosenburg, O. Aharonson, J. W. Head, M. A. Kreslavsky, E. Mazarico, G. A. Neumann, D. E. Smith, M. H. Torrence, and M. T. Zuber, ``Global surface slopes and roughness of the Moon from the Lunar Orbiter Laser Altimeter,'' \emph{Journal of Geophysical Research: Planets}, vol. 116, E02001, 2011. \href{https://doi.org/10.1029/2010JE003716}{doi:10.1029/2010JE003716}.

\bibitem{pgda}
NASA Goddard Space Flight Center, Planetary Geology, Geophysics and Geochemistry Laboratory, ``Lunar slopes and Hurst exponent,'' Product 70. \url{https://pgda.gsfc.nasa.gov/products/70}, accessed 15 Nov 2026.

\bibitem{horizons}
Jet Propulsion Laboratory, Solar System Dynamics, ``Horizons System'' and ``Horizons API'' documentation. \url{https://ssd.jpl.nasa.gov/horizons/} and \url{https://ssd-api.jpl.nasa.gov/doc/horizons.html}, accessed 15 Nov 2026.

\bibitem{smith2010}
D. E. Smith et al., ``The Lunar Orbiter Laser Altimeter Investigation on the Lunar Reconnaissance Orbiter Mission,'' \emph{Space Science Reviews}, vol. 150, pp. 209--241, 2010. \href{https://doi.org/10.1007/s11214-009-9512-y}{doi:10.1007/s11214-009-9512-y}.
\end{thebibliography}

\end{document}
